% Generated draft: requires mathematical and source review.
\documentclass[11pt]{article}
\usepackage[margin=25mm]{geometry}
\usepackage{fontspec}
\setmainfont{Latin Modern Roman}
\usepackage{amsmath,amssymb,mathtools,unicode-math}
\setmathfont{Latin Modern Math}
\usepackage{enumitem,needspace,xurl,hyperref}
\setlength{\parindent}{0pt}
\setlength{\parskip}{5pt}
\setlength{\emergencystretch}{3em}
\newcommand{\N}{\mathbb{N}}
\newcommand{\Z}{\mathbb{Z}}
\newcommand{\Q}{\mathbb{Q}}
\newcommand{\R}{\mathbb{R}}
\newcommand{\C}{\mathbb{C}}
\newcommand{\F}{\mathbb{F}}
\newcommand{\A}{\mathbb{A}}
\newcommand{\PP}{\mathbb{P}}
\newcommand{\E}{\mathbb{E}}
\newcommand{\eps}{\varepsilon}
\newcommand{\dd}{\,\mathrm{d}}
\begin{document}
\section*{UnsolvedMath ID: 800830}
\textbf{Source problem code:} catalogue entry
\subsection{Definitions and mathematical statement}
Let \(\mathcal V\) denote the intended class of vector fields, and let \(E(X)\) denote the energy of \(X\in\mathcal V\). The supplied record does not specify the ambient manifold, geometric structures, regularity, divergence conditions, boundary conditions, or energy functional. These data must be fixed before the problem becomes a fully specified mathematical statement.

The problem asks for a sequence \((H_k)_{k\geq 1}\) of higher helicity invariants with the following property: whenever all invariants preceding the first nonzero invariant of a field vanish, that first nonzero invariant gives a lower bound for the field's energy. A schematic formulation of a nontrivial bound is as follows. Write \(H_k:\mathcal V\to A_k\), where each unspecified value space \(A_k\) has a distinguished zero \(0_k\), and seek functions \(b_k:A_k\setminus\{0_k\}\to(0,\infty)\) such that, for every \(k\geq 1\) and every \(X\in\mathcal V\),
\[
\left[H_j(X)=0_j\text{ for all }1\leq j<k,\quad H_k(X)\neq 0_k\right]
\quad\Longrightarrow\quad
E(X)\geq b_k\bigl(H_k(X)\bigr).
\]
The functions \(b_k\) may depend on the fixed ambient data, but must apply uniformly to the fields in the indicated stratum. For \(k=1\), the preceding-vanishing condition is empty. The requested implication imposes no condition on a field for which every \(H_k\) vanishes.

This display expresses an interpretation of the requested energy bound, not a quantitative formula supplied by the source. The value spaces, meaning of higher helicity, transformation group or equivalence relation underlying invariance, admissible dependence of the bound, and normalization of the energy remain unspecified. The record also does not state whether the sequence must begin with classical helicity.
\subsection{Short English statement}
Construct successive higher helicity invariants so that the first nonzero invariant of a vector field controls its energy from below when every preceding invariant vanishes. The supplied formulation leaves the geometric setting, invariant definitions, and quantitative bound unresolved.
\subsection{Sources}
\begin{itemize}
\item Local UnsolvedMath registry, supplied complete record: permanent ID \(800830\), source problem code \texttt{catalogue entry}, upstream revision \texttt{d371ee917ac9647fbb06cd6fc529deec6f3dda82}. The registry labels the problem open; this supplied status is not an independently verified mathematical or literature assessment.
\item Khesin, Misiolek, Modin, \emph{Geometric Hydrodynamics in Open Problems} (2022), Problem 29. Supplied source URL: \url{https://arxiv.org/abs/2205.01143}. The record reports access on 2026-10-16 and extraction by \texttt{arxiv\_latex}; these are supplied provenance details, not a claim of independent source consultation.
\item The registry lists Khesin, Misiolek, Modin as proposers, with proposed year 2022, and preserves the attribution/context marker \texttt{[ak]}. The supplied record does not expand that marker.
\end{itemize}
\end{document}
